\chapter{Modelo de Memoria y GIL}
\label{chap:memory}
{\huge$\bigstar$}{\huge$\bigstar$} Intermedio

\section{Estructura de VALUE\_PYTHON}

Cada \texttt{VALUE\_PYTHON} envuelve un \texttt{PyObject*} de CPython.
Tres campos: \texttt{obj} (puntero), \texttt{callable} (caché),
\texttt{callable\_known} (flag de caché).

\section{Ownership}

Cada \texttt{VALUE\_PYTHON} posee exactamente una referencia a su
\texttt{PyObject*}. Al destruirse, llama \texttt{Py\_DECREF}.

\section{GIL}

CPython requiere el GIL. LisPy lo gestiona con macros:
\begin{verbatim}
#define PY_GIL_ENSURE()  PyGILState_STATE _py_gil = PyGILState_Ensure()
#define PY_GIL_RELEASE() PyGILState_Release(_py_gil)
\end{verbatim}

El GIL se adquiere por llamada, no por ámbito léxico.

\section{Cachés}

\begin{itemize}
\item \textbf{Módulos:} hash-table \texttt{equal}, persiste toda la vida.
\item \textbf{Callables:} \texttt{PyCallable\_Check} cacheado.
\end{itemize}

